% Response_to_Reviewers.tex
\documentclass[11pt,a4paper]{article}

\usepackage[utf8]{inputenc}
\usepackage[T1]{fontenc}
\usepackage{lmodern}
\usepackage{microtype}

\usepackage{amsmath,amssymb}
\usepackage{geometry}
\usepackage{enumitem}
\usepackage{parskip}
\usepackage{hyperref}
\usepackage{xcolor}

\geometry{
  top=2.5cm,
  bottom=2.5cm,
  left=2.7cm,
  right=2.7cm
}

\hypersetup{
  colorlinks=true,
  linkcolor=black,
  urlcolor=blue,
  citecolor=black
}

\setlength{\parindent}{0pt}
\setlength{\parskip}{0.7em}

% Reviewer-comment formatting
\newenvironment{reviewcomment}
  {\begin{quote}\itshape}
  {\end{quote}}

\newcommand{\response}{\noindent\textbf{Response:} }

\title{\textbf{Response to the Editor and Reviewers}}
\author{}
\date{}

\begin{document}

\maketitle
We thank the editor and both reviewers for their comments. Their observations led us to revise the study and present the results in greater detail. Before changing any numerical values, we audited the code for all four domains and the historical artifacts used to generate the original results. In the revised version, we retain the predictions for each forecasting source and calculate relative capability only after verifying that the comparisons share exact common support. We collect the data needed to reproduce the analysis at a single entry point.


% ============================================================
\section*{Editor}
% ============================================================

\begin{reviewcomment}
Please report results precisely, remove overstated conclusions, and explain
limitations completely.
\end{reviewcomment}

\response
We reconstructed and verified the wind and traffic inputs by reproducing every
submitted persistence and LightGBM horizon error to floating-point tolerance.
In ``Models'' (pp.~10--11), we corrected the fixed LightGBM configurations,
removed unsupported claims of model invariance, and eliminated causal language.
We expanded ``Limitations'' (p.~19) to state explicitly that findings are
conditional upon the chosen baseline, evaluation metric, $H_{\max}$, dataset,
model configuration, and validation protocol. We also state directly that a
dynamic $H^*$ controller was not prospectively evaluated. The new
``Alternative-baseline sensitivity'' subsection (pp.~16--17) and Table~4 report
the additional sensitivity analysis.


% ============================================================
\section*{Reviewer 1}
% ============================================================

\subsection*{Comment 1}

\begin{reviewcomment}
In the operational definitions, use a colon before the explanations of the
relaxed and strict horizons.
\end{reviewcomment}

\response
We replaced the full stop with a colon after both labels in
``Operational Predictability Horizon'' (Section~3.1, p.~6).


\subsection*{Comments 2--3}

\begin{reviewcomment}
Equation (5) and the surrounding text should explicitly state that the relaxed
horizon permits non-positive skill at intermediate steps.
\end{reviewcomment}

\response
We now define
\[
\mathcal{H}^{+}
=
\{h\in\mathcal{H}:\mathrm{Skill}(h)>0\}
\]
and
\[
H^*(\mathrm{relax})
=
\max\left(\mathcal{H}^{+}\cup\{0\}\right)
\]
in Section~3.1 (p.~6).

The text explicitly notes that intermediate horizons with non-positive skill
may fall within the relaxed reach, but are not themselves interpreted as
skilful. We verified this behaviour using executable unit tests covering
intermediate skill deficits and delayed onset.


\subsection*{Comment 4}

\begin{reviewcomment}
Maintain units throughout Table 2.
\end{reviewcomment}

\response
Table~2 (p.~11) now specifies PM$_{2.5}$ in $\mu\mathrm{g\,m^{-3}}$, wind speed in
$\mathrm{m\,s^{-1}}$, and traffic speed in mph. For electric load, we clarify that the
stored experimental target is the sum of 15-minute kW readings and that
division by four expresses the aggregate daily demand in kWh, following the source
metadata and the aggregation script.


\subsection*{Comment 5}

\begin{reviewcomment}
Figures 1--3 are not correctly referenced in the text.
\end{reviewcomment}

\response
We have added explicit textual references to the conceptual $H^*$ schematic
(Fig.~1, p.~7), the rolling-origin validation protocol (Fig.~2, p.~9), and the
empirical skill curve (Fig.~3, p.~12).


\subsection*{Comment 6}

\begin{reviewcomment}
Evaluate baseline choices beyond persistence.
\end{reviewcomment}

\response
Before inspecting any comparative result, we pre-specified seasonal persistence
as a transparent cross-domain sensitivity baseline: 24-hour seasonality for
hourly domains and 7-day seasonality for daily load. LightGBM, persistence, and
seasonal persistence were regenerated across identical forecast origins, and
exact equality of forecast origin, target timestamp, horizon, and
$y_{\mathrm{true}}$ was asserted before Skill was calculated. No records were
dropped during pairwise alignment.

For PM$_{2.5}$, the persistence comparison on this exact common support yields
$H^*(\mathrm{strict})=22$ with longest contiguous positive interval $[27,48]$,
whereas seasonal persistence yields $H^*(\mathrm{strict})=24$ with $[25,48]$.

The matched-support persistence value $H^*(\mathrm{strict})=22$ differs from
the primary full-support persistence result ($H^*(\mathrm{strict})=13$ with
longest contiguous positive interval $[36,48]$) because the sensitivity
analysis is restricted to origins jointly eligible for both baselines. This
difference is an evaluation-support effect: the matched-support persistence row
is used solely for the baseline-sensitivity comparison and does not replace the
primary PM$_{2.5}$ result.

Under the same matched-support analysis, load remains at
$H^*(\mathrm{strict})=1$, wind changes from 48 to 30, and traffic changes from
7 to 14. $H^*(\mathrm{relax})$ remains unchanged across all four domains. We
note that the hourly endpoints are exact multiples of the seasonal period,
where seasonal and standard persistence coincide; the unchanged relaxed
endpoint is therefore partly structural and is not interpreted as evidence of
baseline invariance.


\subsection*{Comment 7}

\begin{reviewcomment}
Explain why ARIMA(2,0,0) was selected and whether ACF was examined.
\end{reviewcomment}

\response
We agree that the original manuscript did not document this choice adequately.
An audit of our repository and Git history revealed no record that the AR(2)
order was selected via ACF/PACF, AIC/BIC, grid search, auto-ARIMA, or validation
splits. We therefore make no such claim.

The audit also revealed that the historical wind adapter requested a one-step
ARIMA forecast at every evaluated horizon, failing to implement the stated
$h$-step protocol. We removed the ARIMA numerical robustness claims from
``Models'' and the Wind and Traffic results sections rather than retaining an
unsupported rationale or modifying the historical experiment post hoc.

The requested robustness evidence is instead provided by the prospectively
specified baseline sensitivity analysis on exact common support
(Section~10.4, pp.~16--17).


\subsection*{Comments 8--9}

\begin{reviewcomment}
Explain the practical utility of relaxed and strict $H^*$ and what they add
beyond fixed-horizon summaries.
\end{reviewcomment}

\response
A new ``Operational interpretation and actionability'' subsection
(Section~3.2, p.~7) defines $H^*(\mathrm{relax})$ as the outer horizon at
which positive baseline-relative skill is observed, and $H^*(\mathrm{strict})$,
with its interval location, as the longest contiguous positive interval.

We explain that fixed-horizon summaries fail to capture delayed onset, skill
deficits, fragmentation, interval placement, continuity, or late recovery.


\subsection*{Comment 10}

\begin{reviewcomment}
Could $H^*$ dynamically decide which horizons to forecast, and could this
create feedback or overfitting?
\end{reviewcomment}

\response
We have added a substantive, bounded discussion to Section~11 (p.~18). In an
operational setting, $H^*$ could inform a decision policy estimated on
historical validation data, frozen for prospective deployment, and subsequently
updated using only realised historical errors.

We explicitly caution that estimating $H^*$, selecting horizons, and
evaluating performance on the same errors introduces adaptive selection bias.
A dynamic controller would require prospective, prequential, or nested
validation, and is not claimed to be validated here.


% ============================================================
\section*{Reviewer 2}
% ============================================================

\subsection*{Comment 1}

\begin{reviewcomment}
Explain the actionability of the findings more explicitly.
\end{reviewcomment}

\response
We added the dedicated operational interpretation subsection noted above
(Section~3.2, p.~7), distinguishing the outer positive reach from the longest
contiguous positive interval and restricting dynamic deployment claims to
future, separately validated policies.


\subsection*{Comment 2}

\begin{reviewcomment}
Figures 1 and 2 are not referenced correctly.
\end{reviewcomment}

\response
Both figures are now explicitly introduced and discussed in the main text
(Fig.~1, p.~7; Fig.~2, p.~9).


\subsection*{Comment 3}

\begin{reviewcomment}
In Figure 1, the curve overlaps text or annotations.
\end{reviewcomment}

\response
We regenerated Figure~1 with dedicated callout boxes and separated label
placements. The underlying data and mathematical definitions remain unchanged,
and the rendered PDF (p.~7) was verified visually.


\subsection*{Comment 4}

\begin{reviewcomment}
In Figure 2, text does not fit inside the boxes.
\end{reviewcomment}

\response
We regenerated Figure~2 with larger containers and wrapped text, resolving the
text overflow and clipping in the rendered PDF (p.~9).


\subsection*{Comment 5}

\begin{reviewcomment}
If possible, evaluate alternative baselines beyond persistence.
\end{reviewcomment}

\response
We conducted the pre-specified seasonal-persistence sensitivity analysis
detailed in our response to Reviewer~1 (Comment~6), enforcing exact common
support and reporting all outcomes regardless of direction. Table~4 reports
all four domains across both baselines, documenting common-origin sample
counts, $H^*(\mathrm{relax})$, $H^*(\mathrm{strict})$, and interval locations.


\subsection*{Comment 6}

\begin{reviewcomment}
Data Availability should reduce the need to visit four or five separate
sources and should clearly state whether files are cached in the code
repository.
\end{reviewcomment}

\response
The replication package now provides a single entry point
(\texttt{data/README.md}) alongside a machine-readable manifest containing
source identifiers, exact filenames, SHA-256 checksums, pre-processing
scripts, and redistribution terms.

Exact input records for PM$_{2.5}$ and load are version-controlled within the
repository. Wind and traffic inputs are retrieved deterministically,
checksum-verified, and pre-processed via a single script; we state explicitly
that raw external files are not cached directly in the repository. The
manuscript directs readers to this manifest in ``Data Availability'' (p.~21).


\subsection*{Comment 7}

\begin{reviewcomment}
Explain the reason for ARIMA(2,0,0).
\end{reviewcomment}

\response
As noted in our response to Reviewer~1 (Comment~7), no documented
order-selection rationale was preserved, and the historical wind implementation
did not execute the intended $h$-step protocol. We removed the ARIMA robustness
claims from the text and document this constraint transparently, rather than
constructing a post hoc ACF/PACF justification.


\end{document}